\documentclass[UTF8,11pt]{ctexart}
\usepackage[a4paper,margin=24mm]{geometry}
\usepackage{amsmath,amssymb,amsthm,mathtools,mathrsfs}
\usepackage[all]{xy}
\usepackage{xurl,hyperref,enumitem}
\usepackage{tikz}
\usetikzlibrary{arrows.meta,cd,decorations.markings,calc,patterns}
\hypersetup{hidelinks,breaklinks=true}
\setlength{\emergencystretch}{3em}
\allowdisplaybreaks
\newtheorem*{theorem}{Theorem}
\newtheorem*{thm}{Theorem}
\newtheorem*{lemma}{Lemma}
\newtheorem*{proposition}{Proposition}
\newtheorem*{prop}{Proposition}
\newtheorem*{corollary}{Corollary}
\newtheorem*{cor}{Corollary}
\newtheorem*{claim}{Claim}
\theoremstyle{definition}
\newtheorem*{definition}{Definition}
\newtheorem*{defn}{Definition}
\newtheorem*{remark}{Remark}
\newtheorem*{example}{Example}
\newcommand{\R}{\mathbb R}
\newcommand{\C}{\mathbb C}
\newcommand{\Z}{\mathbb Z}
\newcommand{\Q}{\mathbb Q}
\newcommand{\N}{\mathbb N}
\newcommand{\F}{\mathbb F}
\newcommand{\D}{\mathbb D}
\newcommand{\bB}{\mathbb B}
\newcommand{\bC}{\mathbb C}
\newcommand{\bR}{\mathbb R}
\newcommand{\bZ}{\mathbb Z}
\newcommand{\bN}{\mathbb N}
\newcommand{\bQ}{\mathbb Q}
\newcommand{\bD}{\mathbb D}
\newcommand{\bF}{\mathbb F}
\newcommand{\CP}{\mathbb{CP}}
\newcommand{\RP}{\mathbb{RP}}
\newcommand{\sO}{\mathscr O}
\newcommand{\sI}{\mathscr I}
\newcommand{\sF}{\mathscr F}
\newcommand{\sG}{\mathscr G}
\newcommand{\sH}{\mathscr H}
\newcommand{\sR}{\mathscr R}
\newcommand{\sM}{\mathscr M}
\newcommand{\sV}{\mathscr V}
\newcommand{\sU}{\mathscr U}
\DeclareMathOperator{\Hom}{Hom}
\DeclareMathOperator{\Ext}{Ext}
\DeclareMathOperator{\Tor}{Tor}
\DeclareMathOperator{\Spec}{Spec}
\DeclareMathOperator{\Proj}{Proj}
\DeclareMathOperator{\supp}{supp}
\DeclareMathOperator{\im}{im}
\DeclareMathOperator{\id}{id}
\DeclareMathOperator{\Tr}{Tr}
\DeclareMathOperator{\tr}{tr}
\DeclareMathOperator{\rank}{rank}
\DeclareMathOperator{\ord}{ord}
\DeclareMathOperator{\dist}{dist}
\DeclareMathOperator{\End}{End}
\DeclareMathOperator{\Aut}{Aut}
\DeclareMathOperator{\Gal}{Gal}
\DeclareMathOperator{\vol}{vol}
\DeclareMathOperator{\Cl}{Cl}
\DeclareMathOperator{\coker}{coker}
\DeclareMathOperator{\codim}{codim}
\DeclareMathOperator{\divergence}{div}
\DeclareMathOperator{\Ric}{Ric}
\DeclareMathOperator{\grad}{grad}
\DeclareMathOperator{\Hess}{Hess}
\newcommand{\abs}[1]{\left\lvert#1\right\rvert}
\newcommand{\sabs}[1]{\lvert#1\rvert}
\newcommand{\babs}[1]{\bigl\lvert#1\bigr\rvert}
\newcommand{\Babs}[1]{\Bigl\lvert#1\Bigr\rvert}
\newcommand{\bbabs}[1]{\biggl\lvert#1\biggr\rvert}
\newcommand{\BBabs}[1]{\Biggl\lvert#1\Biggr\rvert}
\newcommand{\norm}[1]{\left\lVert#1\right\rVert}
\newcommand{\snorm}[1]{\lVert#1\rVert}
\newcommand{\bnorm}[1]{\bigl\lVert#1\bigr\rVert}
\newcommand{\Bnorm}[1]{\Bigl\lVert#1\Bigr\rVert}
\newcommand{\bbnorm}[1]{\biggl\lVert#1\biggr\rVert}
\newcommand{\BBnorm}[1]{\Biggl\lVert#1\Biggr\rVert}
\newcommand{\widebar}[1]{\overline{#1}}
\newcommand{\nicefrac}[2]{\frac{#1}{#2}}
\newcommand{\linnprod}[2]{\langle#1,#2\rangle}
\newcommand{\myindex}[1]{#1}
\newcommand{\glsadd}[1]{}
\newcommand{\avoidbreak}{}
\newcommand{\sourceref}[1]{\textnormal{[原文：\nolinkurl{#1}]}}
\newcommand{\thmref}[1]{\sourceref{#1}}
\newcommand{\lemmaref}[1]{\sourceref{#1}}
\newcommand{\propref}[1]{\sourceref{#1}}
\newcommand{\corref}[1]{\sourceref{#1}}
\newcommand{\defnref}[1]{\sourceref{#1}}
\newcommand{\exampleref}[1]{\sourceref{#1}}
\newcommand{\exerciseref}[1]{\sourceref{#1}}
\newcommand{\figureref}[1]{\sourceref{#1}}
\newcommand{\sectionref}[1]{\sourceref{#1}}
\newcommand{\appendixref}[1]{\sourceref{#1}}
\newcommand{\stacksref}[2]{\href{https://stacks.math.columbia.edu/tag/#1}{#2 [#1]}}

\begin{document}
\section*{037\quad 适当态射的Stein有限分解与几何连通纤维}
\subsection*{来源与计数目标}
The Stacks Project Authors, \emph{The Stacks Project}。Chapter 37, Section 37.53，Theorem 37.53.4；支撑为 Lemmas 37.53.1、37.53.3。使用实际取得的网页数学 TeX 正文，获取日期：2026-09-28；未取得网页构建的固定提交号，不编造版本。
\par\url{https://stacks.math.columbia.edu/tag/03H0}
\par\url{https://stacks.math.columbia.edu/tag/03GX}

\par\textbf{本文件只计一个问题：}局部 Noetherian 基底上的 Stein 分解及几何连通纤维；一般基底版本不另计。
\subsection*{作者命题与人类证明}

\begin{lemma}[37.53.1]
Let $S$ be a scheme. Let $f:X\to S$ be a universally closed and quasi-separated morphism. There exists a factorization
\[\xymatrix{X\ar[rr]_{f'}\ar[rd]_f&&S'\ar[dl]^\pi\\&S&}\]
with the following properties:
\begin{enumerate}
\item the morphism $f'$ is universally closed, quasi-compact, quasi-separated, and surjective,
\item the morphism $\pi:S'\to S$ is integral,
\item we have $f'_*\mathcal O_X=\mathcal O_{S'}$,
\item we have $S'=\underline{\Spec}_S(f_*\mathcal O_X)$, and
\item $S'$ is the normalization of $S$ in $X$, see Morphisms, Definition 29.54.3.
\end{enumerate}
Formation of the factorization $f=\pi\circ f'$ commutes with flat base change.
\end{lemma}
\begin{proof}
By Morphisms, Lemma 29.42.8 the morphism $f$ is quasi-compact. Hence the normalization $S'$ of $S$ in $X$ is defined (Morphisms, Definition 29.54.3) and we have the factorization $X\to S'\to S$. By Morphisms, Lemma 29.54.11 we have (2), (4), and (5). The morphism $f'$ is universally closed by Morphisms, Lemma 29.42.7. It is quasi-compact by Schemes, Lemma 26.21.14 and quasi-separated by Schemes, Lemma 26.21.13.

To show the remaining statements we may assume the base scheme $S$ is affine, say $S=\Spec(R)$. Then $S'=\Spec(A)$ with $A=\Gamma(X,\mathcal O_X)$ an integral $R$-algebra. Thus it is clear that $f'_*\mathcal O_X$ is $\mathcal O_{S'}$ (because $f'_*\mathcal O_X$ is quasi-coherent, by Schemes, Lemma 26.24.1, and hence equal to $\widetilde A$). This proves (3).

Let us show that $f'$ is surjective. As $f'$ is universally closed (see above) the image of $f'$ is a closed subset $V(I)\subset S'=\Spec(A)$. Pick $h\in I$. Then $h|_X=f^\sharp(h)$ is a global section of the structure sheaf of $X$ which vanishes at every point. As $X$ is quasi-compact this means that $h|_X$ is a nilpotent section, i.e., $h^n|X=0$ for some $n>0$. But $A=\Gamma(X,\mathcal O_X)$, hence $h^n=0$. As every element of $I$ is nilpotent, we conclude that $V(I)=S'$ as desired.

By Cohomology of Schemes, Lemma 30.5.2 we see that formation of $f_*\mathcal O_X$ commutes with flat base change. Formation of the relative spectrum commutes with any base change by Constructions, Lemma 27.4.6. Thus formation of the factorization commutes with flat base change.
\end{proof}

\begin{lemma}[37.53.3]
Let $f:X\to S$ be a morphism of schemes. Let $s\in S$. Then $X_s$ is geometrically connected, if and only if for every \'etale neighbourhood $(U,u)\to(S,s)$ the base change $X_U\to U$ has connected fibre $X_u$.
\end{lemma}
\begin{proof}
If $X_s$ is geometrically connected, then any base change of it is connected. On the other hand, suppose that $X_s$ is not geometrically connected. Then by Varieties, Lemma 33.7.11 we see that $X_s\times_{\Spec(\kappa(s))}\Spec(k)$ is disconnected for some finite separable field extension $k/\kappa(s)$. By Lemma 37.35.2 there exists an affine \'etale neighbourhood $(U,u)\to(S,s)$ such that $\kappa(u)/\kappa(s)$ is identified with $k/\kappa(s)$. In this case $X_u$ is disconnected.
\end{proof}

\begin{theorem}[37.53.4; Stein factorization, Noetherian case]
Let $S$ be a locally Noetherian scheme. Let $f:X\to S$ be a proper morphism. There exists a factorization
\[\xymatrix{X\ar[rr]_{f'}\ar[rd]_f&&S'\ar[dl]^\pi\\&S&}\]
with the following properties:
\begin{enumerate}
\item the morphism $f'$ is proper with geometrically connected fibres,
\item the morphism $\pi:S'\to S$ is finite,
\item we have $f'_*\mathcal O_X=\mathcal O_{S'}$,
\item we have $S'=\underline{\Spec}_S(f_*\mathcal O_X)$, and
\item $S'$ is the normalization of $S$ in $X$, see Morphisms, Definition 29.54.3.
\end{enumerate}
\end{theorem}
\begin{proof}
Let $f=\pi\circ f'$ be the factorization of Lemma 37.53.1. Note that besides the conclusions of Lemma 37.53.1 we also have that $f'$ is separated (Schemes, Lemma 26.21.13) and finite type (Morphisms, Lemma 29.15.8). Hence $f'$ is proper. By Cohomology of Schemes, Proposition 30.19.1 we see that $f_*\mathcal O_X$ is a coherent $\mathcal O_S$-module. Hence we see that $\pi$ is finite, i.e., (2) holds.

This proves all but the most interesting assertion, namely that all the fibres of $f'$ are geometrically connected. It is clear from the discussion above that we may replace $S$ by $S'$, and we may therefore assume that $S$ is Noetherian, affine, $f:X\to S$ is proper, and $f_*\mathcal O_X=\mathcal O_S$. Let $s\in S$ be a point of $S$. We have to show that $X_s$ is geometrically connected.

By Lemma 37.53.3 we see that it suffices to show $X_u$ is connected for every \'etale neighbourhood $(U,u)\to(S,s)$. We may assume $U$ is affine. Thus $U$ is Noetherian (Morphisms, Lemma 29.15.6), the base change $f_U:X_U\to U$ is proper (Morphisms, Lemma 29.42.5), and that also $(f_U)_*\mathcal O_{X_U}=\mathcal O_U$ (Cohomology of Schemes, Lemma 30.5.2).

Hence after replacing $(f:X\to S,s)$ by the base change $(f_U:X_U\to U,u)$ it suffices to prove that the fibre $X_s$ is connected when $f_*\mathcal O_X=\mathcal O_S$. We can deduce this from Derived Categories of Schemes, Lemma 36.32.7 (by looking at idempotents in the structure sheaf of $X_s$) but we will also give a direct argument below.

Namely, we apply the theorem on formal functions, more precisely Cohomology of Schemes, Lemma 30.20.7. It tells us that
\[\mathcal O_{S,s}^{\wedge}=(f_*\mathcal O_X)_s^{\wedge}=\varprojlim_n H^0(X_n,\mathcal O_{X_n})\]
where $X_n$ is the $n$th infinitesimal neighbourhood of $X_s$. Since the underlying topological space of $X_n$ is equal to that of $X_s$ we see that if $X_s=T_1\amalg T_2$ is a disjoint union of nonempty open and closed subschemes, then similarly $X_n=T_{1,n}\amalg T_{2,n}$ for all $n$. And this in turn means $H^0(X_n,\mathcal O_{X_n})$ contains a nontrivial idempotent $e_{1,n}$, namely the function which is identically $1$ on $T_{1,n}$ and identically $0$ on $T_{2,n}$.

It is clear that $e_{1,n+1}$ restricts to $e_{1,n}$ on $X_n$. Hence $e_1=\varprojlim e_{1,n}$ is a nontrivial idempotent of the limit. This contradicts the fact that $\mathcal O_{S,s}^{\wedge}$ is a local ring. Thus the assumption was wrong, i.e., $X_s$ is connected, and we win.
\end{proof}

\subsection*{整理说明（非作者正文）}
主定理的分解构造、满射性、\'etale 基变换归约和幂等元矛盾均已收录。明确引用为先修的结果包括相对谱与相对正规化、适当推送保持凝聚性、平坦基变换、形式函数定理，以及完备局部环仍是局部环。并未将“结构层推送等于基底即纤维连通”当作先修，该核心论证以无穷小邻域上的相容幂等元展开。

作者为连通性同时给出一个导出范畴引理引用和随后直接证明，两者的关系及直接证明均保留。保留原文的 $f^\sharp(h)$ 写法与论证次序；极限符号只作 TeX 规范化。难度来源是通过形式邻域的函数代数识别几何断开，再利用局部代数性质排除该断开，不只是对现成 Stein 分解作应用。
\subsection*{可修改的简短标注（非作者正文）}
$c$：适当态射的有限分解与纤维连通性。

$a$：相对谱；相对正规化；凝聚推送；平坦基变换；形式函数定理；无穷小邻域；幂等元；完备局部环。

\texttt{cross\_theory\_status = yes}。几何纤维的断开被转成相容幂等元，经形式函数定理进入完备局部环并得到矛盾，因而确定分解的几何连通性质。位置：Theorem 37.53.4 的后半证明。

全部标注均为非穷尽、可修改初稿；未标注不表示未使用。
\subsection*{许可与修改记录}
原数学正文作者：The Stacks Project Authors。按 GNU Free Documentation License 1.2 或后续版本复用；无不变章节、封面文字或封底文字。许可全文位于 \texttt{sources/stacks/GFDL-1.2.txt}。本条整理部分沿用同一许可。2026-09-28：助手摘录、独立排版并加入来源、范围说明与初稿标注；原作者不对本次整理背书。
\end{document}
